\documentclass[11pt]{article}
\usepackage[a4paper,margin=1in]{geometry}
\usepackage{amsmath,amssymb,amsthm,mathtools}
\usepackage[hidelinks]{hyperref}
\usepackage{booktabs}

\newtheorem{theorem}{Theorem}
\newtheorem{proposition}[theorem]{Proposition}
\newtheorem{lemma}[theorem]{Lemma}
\newcommand{\Cl}{\operatorname{Cl}}
\newcommand{\Tr}{\operatorname{Tr}}
\newcommand{\Q}{\mathbb Q}
\newcommand{\Z}{\mathbb Z}

\title{Twisted-Convolution Identities in Dimensions Four and Five\\
from Shintani Ray Units}
\author{Hainan Zhao\\
\small Independent Researcher\\
\small\texttt{hainzhao@gmail.com}}
\date{28 July 2026}

\begin{document}
\maketitle

\begin{abstract}
We prove the formal Twisted Convolution Conjecture unconditionally in
dimensions four and five.  Dimension four is the quadratic calibration:
the principal-ghost table depends on one positive double-sine value,
all $36$ rank-two minors have one common quadratic factor, and Kopp's
Kronecker-limit formula together with an analytic class-number
calculation evaluates that factor exactly.  Dimension five is the main
arithmetic case.  For $K=\Q(\sqrt3)$ and modulus $(5)\infty_2$, we
compute all twenty-four nonzero characteristics, their positive Kopp
lifts, ray classes, signs, and multipliers.  Shintani's 1978
algebraicity theorem gives an explicit safe exponent $5760$ that makes
his existential algebraicity effective; labeled ray-field
and Frobenius certificates, rigorous Arb enclosures, and Voutier's
height bound then identify the complete reciprocal degree-eight ray
packet.  Exact number-field arithmetic makes all $100$ rank-two
minors vanish.  In both dimensions complex conjugation supplies the
paired shift, and integral form covariance gives the result for every
admissible tuple.  No Stark conjecture or \texttt{bnrstark}
recognition is used as an input.  We finish with exact research-status
boundaries for dimensions six and eight: one primitive order-six
Shintani--Stark value remains in dimension six, while two oriented
cyclic-quartic resolvents remain in dimension eight.  For the latter
we identify an exact projective-$V_4$ descent gate to the imaginary
quadratic fields $\Q(\sqrt{-6})$ and $\Q(\sqrt{-30})$.  The finite
cancellation certificates in those dimensions are already exact, but
the boundary results are not asserted as unconditional TCC theorems.
\end{abstract}

\section{Overview and common convention}

The principal-ghost construction of Appleby--Flammia--Kopp
\cite{AFK}, continuing the SIC--Stark program of Kopp
\cite{Kopp2018}, associates convention-sensitive Shintani--Faddeev
real-multiplication values to a finite Weyl reconstruction.  The formal
TCC asks for two shifted twisted convolutions to vanish, equivalently
for the reconstructed ghost matrices to be rank-one idempotents.  Our
proof architecture is
\[
\begin{array}{c}
\text{characteristic-to-ray and multiplier bridge}\\
\Downarrow\\
\text{analytic identification of a labeled ray-unit packet}\\
\Downarrow\\
\text{exact vanishing of every rank-two minor}\\
\Downarrow\\
\text{both formal shifts and integral-form transport}.
\end{array}
\]

To remove the standard double-sine ambiguity once and for all, define
\[
 S_2(z\mid\omega_1,\omega_2)
 :=\frac{\Gamma_2(z\mid\omega_1,\omega_2)}
 {\Gamma_2(\omega_1+\omega_2-z\mid\omega_1,\omega_2)}.
\tag{1}
\]
Thus $S_2$ is the reciprocal of Kopp's
$\operatorname{Sin}_{2,\mathrm{Kopp}}$ convention
\cite[Definition~4.22, equation~(4.8)]{Kopp}.

\subsection{AFK terminology and the formal TCC}

We recall the data used in the theorem; this also fixes our transcription
of AFK equation~(1.49).  Following AFK Definitions~1.21--1.24, a pair
$(d,r)$ is \emph{admissible} when
\[
 0<r<\frac{d-1}{2},\qquad nr(d-r)=d^2-1
\]
for an integer $n>4$.  Its associated real quadratic field is
$K=\Q(\sqrt{n(n-4)})$.  If $\epsilon>1$ is the fundamental totally
positive unit of $K$, $\Delta_0=\operatorname{disc}(K)$, and
\[
 f_j=\frac{\epsilon^j-\epsilon^{-j}}{\sqrt{\Delta_0}},\qquad
 r_{j,m}=\frac{f_{jm}}{f_j},\qquad
 d_{j,m}=r_{j,m+1}+r_{j,m},\qquad d_j=d_{j,1},
\]
then $(d,r)$ is equivalently indexed by $(K,j,m)$ with
$d=d_{j,m}$ and $r=r_{j,m}$.  An admissible tuple
\[
 t=(d,r,Q)\sim(K,j,m,Q)
\]
also contains a primitive integral indefinite form $Q$ whose fundamental
discriminant is $\Delta_0$ and whose conductor $f$ divides $f_j$
\cite[Definitions~1.27--1.28]{AFK}.
For $Q=\langle a,b,c\rangle$ we use the same symbol for the scaled
Hessian $\left(\begin{smallmatrix}a&b/2\\b/2&c\end{smallmatrix}\right)$,
as AFK do, and write
$\langle(p_1,p_2),(q_1,q_2)\rangle=p_2q_1-p_1q_2$ and
$\omega_d=e^{2\pi i/d}$.

Let $A_t$ be AFK's signed generator of the level-$d$ stability group
of $Q$, let
\[
 S=\begin{pmatrix}0&-1\\1&0\end{pmatrix},\qquad
 L_{z,t}=\frac{d_j-1}{2}I+\frac{f_j}{f}SQ,
\]
and put
\[
 s_d(\boldsymbol p)
 =d+(1+d)(1+p_1)(1+p_2).
\]
For $\boldsymbol p\in\Z^2$, let
$\mathcal I_{\boldsymbol p}$ be any complete set of representatives
for $\Z^2/d\Z^2$ containing $\boldsymbol0$ and $\boldsymbol p$.
AFK Definition~1.34 calls $\lambda\in\Z/d\Z$ a \emph{shift} for $t$
when $2\lambda+d_j-1$ is coprime to $d$ and, for every
$\boldsymbol p$,
\[
 \sum_{\boldsymbol q\in\mathcal I_{\boldsymbol p}}
 \omega_d^{\,r\langle\boldsymbol p,
  (\lambda I+L_{z,t})\boldsymbol q\rangle}
 \operatorname{shin}_{A_t}^{\boldsymbol q/d}(\rho_t)
 \operatorname{shin}_{A_t^{-1}}^{
       (\boldsymbol q-\boldsymbol p)/d}(\rho_t)
 =d^2\delta_{\boldsymbol p,\boldsymbol0}^{(d)}.
\tag{TCC}
\]
Here $\rho_t$ is the positive root of $Q$, and the phases implicit in
$\operatorname{shin}$ use the sign $(-1)^{s_d(\boldsymbol p)}$ from
AFK Definition~1.30.  The set of shifts is denoted $\mathcal Z_t$.
The \emph{formal TCC}, AFK Conjecture~1.35, asserts
$0,1\in\mathcal Z_t$ for every admissible tuple and equality of shift
sets for admissible forms of the same discriminant.  The twist attached
to a shift satisfies
\[
 \det(G)\,f_t(\lambda)\equiv1\pmod{\bar d},\qquad
 f_t(\lambda):=r(2\lambda+d_j-1+d),
\]
where $\bar d=d$ for odd $d$ and $\bar d=2d$ for even $d$.

\begin{theorem}[Low-dimensional TCC]\label{thm:lowdim}
For every admissible tuple in dimension four or five,
\[
 0,1\in\mathcal Z_t.
\]
If $t,t'$ have the same dimension, then
$\mathcal Z_t=\mathcal Z_{t'}$.
\end{theorem}

The dimension-four proof below is deliberately compact: it records the
complete analytic and finite bridge but leaves the thirty-six explicit
minor quotients to the supplementary certificate.  Dimension five
requires the longer powered-algebraicity and height-rigidity argument
that forms the body of the paper.

\section{Dimension four: the quadratic calibration}\label{sec:d4}

Put
\[
 \phi=\frac{1+\sqrt5}{2},\qquad
 \beta_4=\phi^2,\qquad
 Q_4=\langle1,-3,1\rangle,\qquad
 \tau_4=-e^{\pi i/4}.
\]
The complete signed reconstruction table is
\[
T_4(x)=
\begin{pmatrix}
\sqrt5&-x&1&-x^{-1}\\
-x^{-1}&-x^{-1}&-x^{-1}&-x\\
1&-x^{-1}&1&x\\
-x&-x^{-1}&x&-x
\end{pmatrix},
\tag{D4.1}
\]
where
\[
x=\sqrt2\,
\frac{S_2(\beta_4/4\mid\beta_4,1)
      S_2(1/4\mid\beta_4,1)}
     {S_2((\beta_4+1)/4\mid\beta_4,1)}.
\tag{D4.2}
\]
The factor $\sqrt2$ is the reciprocal shift multiplier
$2\sin(\pi/4)$ obtained from \textup{(D4.FE)} at $z=\beta_4/4$.
\begin{lemma}[Dimension-four reconstruction table]\label{lem:d4overlaps}
For the canonical tuple, the fifteen nonzero AFK analytic overlaps are
the nonzero-characteristic entries of \textup{(D4.1)}.  The entry
$(T_4)_{0,0}=\sqrt5$ packages the normalized ghost identity coefficient
into the same table.
\end{lemma}

\begin{proof}
Put
\[
 \mathcal D^{(4)}_{c,b}:=
 S_2\!\left(\frac{c+b\beta_4}{4}\,\middle|\,\beta_4,1\right).
\]
For $(p,q)\ne(0,0)$ and $r\equiv-p-q\pmod4$, the principal-form
specialization of AFK Definition~1.32 is
\[
 \widetilde\nu_{p,q}
 =(-1)^{\,4(p+q)+pq+\min(4,p+q)}
 \mathcal D^{(4)}_{4-p,q}\mathcal D^{(4)}_{4-r,p}
 \mathcal D^{(4)}_{4-q,r}.
\tag{D4.DS}
\]
This formula already includes the AFK phase.  The complete substitution
audit is
\[
\begin{array}{c|c|c@{\qquad}c|c|c}
(p,q)&r&\widetilde\nu_{p,q}&(p,q)&r&\widetilde\nu_{p,q}\\ \hline
(0,1)&3&-x &(0,2)&2&1\\
(0,3)&1&-x^{-1}&(1,0)&3&-x^{-1}\\
(1,1)&2&-x^{-1}&(1,2)&1&-x^{-1}\\
(1,3)&0&-x&(2,0)&2&1\\
(2,1)&1&-x^{-1}&(2,2)&0&1\\
(2,3)&3&x&(3,0)&1&-x\\
(3,1)&0&-x^{-1}&(3,2)&3&x\\
(3,3)&2&-x&&&
\end{array}
\]
Every nonzero-characteristic entry follows from \textup{(D4.DS)}
using only
\[
 S_2(z)S_2(\beta_4+1-z)=1,\quad
 S_2(z+1)=2\sin(\pi z/\beta_4)S_2(z),\quad
 S_2(z+\beta_4)=2\sin(\pi z)S_2(z),
\tag{D4.FE}
\]
and $\beta_4^2=3\beta_4-1$.  The value $\sqrt5$ at $(0,0)$ is not an
analytic zero-characteristic cocycle value: after division by
$\sqrt5$ it represents the Weyl coefficient $\mu_{\boldsymbol0}=1$.
\end{proof}

Define
\[
 K^{(4)}_{r,c}(x)=\frac1{4\sqrt5}
 \sum_{\substack{0\le a,b<4\\r\equiv c+a\pmod4}}
 (T_4)_{a,b}(x)\tau_4^{ab}i^{bc}.
\tag{D4.3}
\]
The representative reduction in (D4.3) retains the even-dimensional
wrap sign.  Its diagonal Fourier sum gives
\[
 \Tr K^{(4)}=(T_4)_{0,0}/\sqrt5=1.
\tag{D4.4}
\]
Exact Laurent reduction proves that every one of its $36$ rank-two
minors belongs to the principal ideal
\[
 \left(x^2-\sqrt{3+\sqrt5}\,x+1\right)
 \subset\Q(i,\sqrt2,\sqrt5)[x,x^{-1}].
\tag{D4.5}
\]
The certificate records all $36$ quotients coefficientwise in the
genuine basis $1,\sqrt2,\sqrt5,\sqrt{10}$; two minors are structurally
zero and the other $34$ have nonzero formal quotients.

It remains to evaluate (D4.2).  Write $K_4=\Q(\sqrt5)$,
$\mathcal O_{K_4}=\Z[\phi]$, and label its embeddings by
$\infty_1(\sqrt5)=+\sqrt5$ and $\infty_2(\sqrt5)=-\sqrt5$.
Since $2$ is inert,
$|(\mathcal O_{K_4}/4)^\times|=12$.  The images of
$\{\pm\phi^n\}$ together with the sign at $\infty_2$ have order $12$;
the ray exact sequence and $h_{K_4}=1$ therefore give
\[
 \Cl_{(4)\infty_2}(K_4)\cong C_2.
\tag{D4.6}
\]
The field
\[
 L_4=K_4(\sqrt\phi)
\]
has relative discriminant $(4)$, and $\infty_2$ ramifies because
$\phi'<0$.  The quadratic conductor--discriminant theorem and
\textup{(D4.6)} show that $L_4$ is the full ray field.

Set $u=\phi+\sqrt\phi$ and let
$\sigma(\sqrt\phi)=-\sqrt\phi$.  Then
$u^\sigma=u^{-1}$ and
\[
 N_{L_4/K_4}(u)=1,\qquad \Tr_{L_4/K_4}(u)=1+\sqrt5.
\]
The field $L_4$ has signature $(2,1)$ and discriminant $-400$.
For $s=\sqrt\phi$, the power-basis discriminant of
$1,s,s^2,s^3$ is also $-400$, equal to the certified field
discriminant; hence $\mathcal O_{L_4}=\Z[s]$.  Its Minkowski bound is
$15/(2\pi)<2.4$, and there is no ideal of norm two because
\[
 X^4-X^2-1\equiv(X^2+X+1)^2\pmod2.
\]
Thus $h_{L_4}=1$.  PARI/GP~2.15.4~\cite{PARI} returns the fundamental units
$s^3-s=s^{-1}$ and $-s^2+s=-u^{-1}$; its
\texttt{bnfcertify}=1 output certifies unconditionally that
\[
 R_{L_4}=\log\phi\,\log u.
\]
For the nontrivial quadratic ray character $\chi$,
$L(s,\chi)=\zeta_{L_4}(s)/\zeta_{K_4}(s)$.  Using
\[
 \zeta_F(s)=-\frac{h_FR_F}{w_F}s^{r_1+r_2-1}
 +O(s^{r_1+r_2})
\]
for $F=K_4,L_4$ gives
\[
 L'(0,\chi)=\log u.
\tag{D4.7}
\]

We now spell out the convention-sensitive specialization of Kopp's
Theorem~1.1~\cite{Kopp}.  Take modulus $4\mathcal O_{K_4}$, identity ray class
$\mathfrak A$, $\mathfrak b=\mathcal O_{K_4}$, $\alpha=4$,
characteristic
\[
 \boldsymbol r=\binom0{1/4},
\qquad
 B_4=
 \begin{pmatrix}3&-1\\1&0\end{pmatrix}^{\!3}
 =
 \begin{pmatrix}21&-8\\8&-3\end{pmatrix}.
\]
Here
$4(\beta_4\Z+\Z)=\mathfrak b\mathfrak m$,
$4(r_2\beta_4-r_1)\mathcal O_{K_4}
=\beta_4\mathcal O_{K_4}=\mathcal O_{K_4}\in\mathfrak A$, and
$r_2\beta_4'-r_1=\beta_4'/4>0$.
Moreover
$B_4\boldsymbol r-\boldsymbol r=(-2,-1)^{\mathsf T}$;
$B_4$ is the positive generator of $\Gamma_{\boldsymbol r}$ with
expanding eigenvalue $\beta_4^3$.  The fiber of
\[
 \Cl_{(4)\infty_1\infty_2}(K_4)
 \longrightarrow \Cl_{(4)\infty_2}(K_4)
\]
has order two, so Kopp's exponent is $n=2/2=1$.

Kopp defines a sign class $\mathfrak R$ in
\cite[Definition~6.3]{Kopp}.  Here it is represented by
$(-1\bmod4,+\text{ at }\infty_2)$.  The residues of
$\phi^n=a_n+b_n\phi$ for $0\le n<6$ are
\[
 (a_n,b_n)=(1,0),(0,1),(1,1),(1,2),(2,3),(3,1),
\]
and $\phi^6\equiv1\pmod4$.  Hence every unit congruent to $-1$ modulo
$4$ is $-\phi^{6k}$, which is negative at $\infty_2$.  Thus
$\mathfrak R\ne\mathfrak A$; by \textup{(D4.6)} it is the nonidentity
class, and the differenced partial zeta is
\[
 Z_{(4)\infty_2}(s,\mathfrak A)
 =\zeta(s,\mathfrak A)-\zeta(s,\mathfrak R\mathfrak A)
 =L(s,\chi).
\tag{D4.8}
\]

The Rademacher invariant of $B_4$ is zero:
$s(21,8)=-1/16$ gives
$18/8-3-12(-1/16)=0$.  Thus $\psi^2(B_4)=1$, while Kopp's theta
character gives
\[
 \chi_{\boldsymbol r}(B_4)
 =\exp\!\left(2\pi i\,\frac12
 \left(-\frac{28}{4}+\frac{152}{16}\right)\right)=i.
\]
Consequently
\[
 (\psi^{-2}\chi_{\boldsymbol r}^{-1})(B_4)=-i
 =\Phi_{(0,1)}^2,\qquad
 \Phi_{(0,1)}=e^{3\pi i/4}.
\tag{D4.9}
\]

For completeness, put
$C=\left(\begin{smallmatrix}3&-1\\1&0\end{smallmatrix}\right)$.
The characteristic orbit in the three cocycle factors is
\[
 (0,1)\longmapsto(3,0)\longmapsto(1,3)\longmapsto(0,1),
\]
with unreduced steps
\[
 C(0,1)=(-1,0)=(3,0)+(-4,0),\quad
 C(3,0)=(9,3)=(1,3)+(8,0),\quad
 C(1,3)=(0,1).
\]
Substitution in \textup{(D4.DS)} gives exponent
$4(0+1)+0+\min(4,1)=5$ and the three arguments
$1+\beta_4/4$, $1/4$, and $1+(3\beta_4-1)/4$.  Hence
\[
 \nu_{(0,1)}
 =\Phi_{(0,1)}
  \operatorname{shin}^{\boldsymbol r}_{B_4}(\beta_4)=-x.
\tag{D4.10}
\]
Squaring \textup{(D4.10)} and applying Kopp's formula line by line now
gives
\[
\begin{split}
 x^2
 &=\Phi_{(0,1)}^2
   \operatorname{shin}^{\boldsymbol r}_{B_4}(\beta_4)^2\\
 &=(\psi^{-2}\chi_{\boldsymbol r}^{-1})(B_4)
   \operatorname{shin}^{\boldsymbol r}_{B_4}(\beta_4)^2\\
 &=\exp Z_{(4)\infty_2}'(0,\mathfrak A)
  =\exp L'(0,\chi)=u.
\end{split}
\tag{D4.11}
\]
Since $u+u^{-1}=1+\sqrt5$ and $x>0$,
\[
 x+x^{-1}=\sqrt{3+\sqrt5}.
\]
Thus (D4.5) makes every minor vanish; (D4.4) then gives
\[
 \operatorname{rank}K^{(4)}=1,\qquad
 (K^{(4)})^2=K^{(4)}.
\tag{D4.12}
\]

\begin{lemma}[Dimension-four coefficient bridge]\label{lem:d4tcc}
For the canonical tuple
$t_4=(4,1,\langle1,-3,1\rangle)$, the equality
$(K^{(4)})^2=K^{(4)}$ is equivalent to the $\lambda=1$ instance of
AFK equation~(1.49).
\end{lemma}

\begin{proof}
For $\boldsymbol p\ne\boldsymbol0$, let
$\nu_{\boldsymbol p}$ denote the analytic value from
Lemma~\ref{lem:d4overlaps}.  Reciprocity gives
\[
 \nu_{\boldsymbol p}\nu_{-\boldsymbol p}=1.
\tag{D4.13}
\]
The actual ghost reconstruction is
\[
 K^{(4)}=\frac14I+\frac1{4\sqrt5}
 \sum_{\boldsymbol p\ne\boldsymbol0}
 \nu_{\boldsymbol p}D_{\boldsymbol p};
\tag{D4.14}
\]
the artificial table entry $(T_4)_{0,0}=\sqrt5$ packages the first
term and is not used as an analytic cocycle value.

Expand \textup{(D4.14)} using
$D_{\boldsymbol p}D_{\boldsymbol q}
=\tau_4^{\langle\boldsymbol p,\boldsymbol q\rangle}
D_{\boldsymbol p+\boldsymbol q}$.  After multiplication by $80$, the
coefficient of $D_{\boldsymbol p}$ for
$\boldsymbol p\ne\boldsymbol0$ is
\[
 \sum_{\substack{\boldsymbol q\bmod4\\
       \boldsymbol q\ne\boldsymbol0,\boldsymbol p}}
 \tau_4^{\langle\boldsymbol p,\boldsymbol q\rangle}
 \nu_{\boldsymbol p-\boldsymbol q}\nu_{\boldsymbol q}
 -2\sqrt5\,\nu_{\boldsymbol p}.
\tag{D4.15}
\]
By \textup{(D4.13)}, the product is
$\nu_{\boldsymbol q}/\nu_{\boldsymbol q-\boldsymbol p}$.

AFK's $\nu_{\boldsymbol0}^{\rm aux}$ is a separate auxiliary cocycle
value.  Their zero-characteristic lemma gives
\[
 \nu_{\boldsymbol0}^{\rm aux}
 +(\nu_{\boldsymbol0}^{\rm aux})^{-1}
 =-(4-2)\sqrt5=-2\sqrt5.
\tag{D4.16}
\]
Thus adjoining the omitted endpoints
$\boldsymbol q=\boldsymbol0,\boldsymbol p$ to \textup{(D4.15)}
produces exactly the full endpoint-corrected sum
\[
 \sum_{\boldsymbol q\in\mathcal I_{\boldsymbol p}}
 \tau_4^{\langle\boldsymbol p,\boldsymbol q\rangle}
 \frac{\nu_{\boldsymbol q}}
 {\nu_{\boldsymbol q-\boldsymbol p}},
\tag{D4.17}
\]
with the endpoint ratios interpreted through
$\nu_{\boldsymbol0}^{\rm aux}$.

For nonzero characteristics write
\[
 \nu_{\boldsymbol q}
 =\Phi_{\boldsymbol q}
 \operatorname{shin}_{A_{t_4}}^{\boldsymbol q/4}(\beta_4),
 \qquad A_{t_4}=B_4.
\]
The two phase identities in AFK Section~5.4 give
\[
 \frac{\Phi_{\boldsymbol q}}{\Phi_{\boldsymbol q-\boldsymbol p}}
 =(-1)^{1+s_4(\boldsymbol p)}\tau_4^{Q_4(\boldsymbol p)}
  \tau_4^{(4-3)\langle\boldsymbol p,\boldsymbol q\rangle}
  \omega_4^{\langle\boldsymbol p,L_{z,t_4}\boldsymbol q\rangle},
\]
\[
 \frac{\operatorname{shin}_{A_{t_4}}^{\boldsymbol q/4}(\beta_4)}
 {\operatorname{shin}_{A_{t_4}}^{(\boldsymbol q-\boldsymbol p)/4}
  (\beta_4)}
 =
 \operatorname{shin}_{A_{t_4}}^{\boldsymbol q/4}(\beta_4)
 \operatorname{shin}_{A_{t_4}^{-1}}^{
   (\boldsymbol q-\boldsymbol p)/4}(\beta_4).
\]
For the identity twist $f_{t_4}(1)=1$.  Its Weyl chirp contributes
the additional
$\tau_4^{\langle\boldsymbol p,\boldsymbol q\rangle}$, and hence
\[
 \tau_4^{(1+4-3)\langle\boldsymbol p,\boldsymbol q\rangle}
 =\omega_4^{\langle\boldsymbol p,\boldsymbol q\rangle}.
\]
Up to a nonzero factor independent of $\boldsymbol q$, the coefficient
equation is therefore
\[
 \sum_{\boldsymbol q\in\mathcal I_{\boldsymbol p}}
 \omega_4^{\langle\boldsymbol p,
 (I+L_{z,t_4})\boldsymbol q\rangle}
 \operatorname{shin}_{A_{t_4}}^{\boldsymbol q/4}(\beta_4)
 \operatorname{shin}_{A_{t_4}^{-1}}^{
       (\boldsymbol q-\boldsymbol p)/4}(\beta_4)
 =16\delta_{\boldsymbol p,\boldsymbol0}^{(4)}.
\tag{D4.TCC}
\]
This is exactly AFK equation~(1.49) with $r=1$ and $\lambda=1$.
For $\boldsymbol p=\boldsymbol0$, the cocycle inverse identity makes
the sum $16$ automatically.  Every step is reversible, including the
endpoint correction \textup{(D4.16)}.
\end{proof}

By \textup{(D4.4)} and \textup{(D4.12)}, $K^{(4)}$ is a trace-one
rank-one idempotent, so Lemma~\ref{lem:d4tcc} proves that $1$ is a
shift.  The twist--shift congruence for the identity twist is
\[
 2\lambda+7\equiv1\pmod8,
\]
hence $\lambda\equiv1\pmod4$.

Complex conjugation replaces the twist $G$ by $J_0G$, where
$J_0=\operatorname{diag}(1,-1)$.  It preserves idempotence and sends
\[
 \lambda\longmapsto-(\lambda+d_j-1)=1-\lambda\pmod4.
\]
It therefore supplies the $\lambda=0$ equation.  Equivalently,
$\det(J_0G)=-1$ gives
$-(2\lambda+7)\equiv1\pmod8$ and hence
$\lambda\equiv0\pmod4$.

Finally, dimension-four admissibility forces
$r=j=m=1$, $K_4=\Q(\sqrt5)$, and conductor one, so every admissible
form has discriminant five.  Class number one and the norm-$-1$ unit
$\phi$ give narrow class number one; all such forms lie in one proper
equivalence class.  The AFK Section~7.1 transformation with
$\boldsymbol q\mapsto\ell M\boldsymbol q$ gives
$\mathcal Z_t\subseteq\mathcal Z_{t_M}$, and applying it to $M^{-1}$
gives the reverse inclusion.  Thus both shifts hold for every
dimension-four admissible tuple and their shift sets are equal.

\section{Dimension five: canonical data and conventions}

Put
\[
 K=\Q(\sqrt3),\qquad \beta=2+\sqrt3,\qquad
 Q=\langle1,-4,1\rangle,\qquad
 L=\begin{pmatrix}4&-1\\1&0\end{pmatrix}.
\]
Then
\[
 B=L^3=\begin{pmatrix}56&-15\\15&-4\end{pmatrix},\qquad
 B\binom0{1/5}-\binom0{1/5}=\binom{-3}{-1}.
\]
Label the real embeddings by
$\infty_1(\sqrt3)=+\sqrt3$ and
$\infty_2(\sqrt3)=-\sqrt3$.  The modulus throughout is
\[
 \mathfrak m=(5)\infty_2.
\]

We continue to use the reciprocal double-sine convention (1).
For $0\le c\le5$ and $0\le b<5$, write
\[
 \mathcal D_{c,b}:=
 S_2\!\left(\frac{c+b\beta}{5}\,\middle|\,\beta,1\right).
\]
For $(p,q)\ne(0,0)$, let $r\equiv-p-q\pmod5$.  The principal-form
specialization of the AFK three-factor formula is
\[
 \widetilde\nu_{p,q}
 =(-1)^{5(p+q)+pq+\min(5,p+q)}
 \mathcal D_{5-p,q}\mathcal D_{5-r,p}\mathcal D_{5-q,r}.
\tag{2}
\]
Formula (2) is used only for nonzero characteristics.  Later we set the
table entry $T_{0,0}:=\sqrt6$ solely so that the identity Weyl
coefficient is $T_{0,0}/\sqrt6=1$.  This reconstruction convention is
not AFK's auxiliary zero-characteristic cocycle value.

\begin{lemma}[Derivation and convention check for (2)]
Formula \textup{(2)} is the specialization of AFK Definition~1.32 to
$t=(5,1,\langle1,-4,1\rangle)$, with no unrecorded reciprocal or sign.
\end{lemma}

\begin{proof}
AFK Definition~1.32 first forms
$\widetilde\nu_{\boldsymbol p}
=\Phi_{\boldsymbol p}\operatorname{shin}^{\boldsymbol p/5}_{B}(\beta)$.
Insert the defining modular-to-Jacobi quotient
\cite[equation (1.26)]{AFK}, and then use the Jacobi cocycle
decomposition \cite[equations (8.2)--(8.5)]{AFK} and the $S$-kernel
formula \cite[equation (8.6)]{AFK}.  In the notation of AFK Section~8.2,
\[
 L=T^4S,\qquad B=L^3,\qquad L\beta=\beta,\qquad
 j_{L^k}(\beta)=\beta^k .
\]
Thus the word $B=(T^4S)^3$ contributes exactly three $S$-kernels.
Modulo five the characteristic orbit is
\[
 (p,q)\longmapsto(r,p)\longmapsto(q,r)\longmapsto(p,q),
 \qquad r\equiv-p-q\pmod5,
\]
and the three unreduced linear arguments are
\[
 \frac{q\beta-p}{5},\qquad
 \frac{p\beta-r}{5},\qquad
 \frac{r\beta-q}{5}.
\]
Reducing their first coordinates to the representatives
$5-p,5-r,5-q$ gives respectively
$\mathcal D_{5-p,q}$, $\mathcal D_{5-r,p}$, and
$\mathcal D_{5-q,r}$.  AFK's $S$-kernel is its displayed double sine
in the denominator (AFK Section~8.2, equation for
$\sigma_S$).  That double sine is Kopp's
$\operatorname{Sin}_{2,\mathrm{Kopp}}$; hence each denominator is
precisely one factor of the reciprocal convention (1).

For completeness, apply the integral-translation law
\cite[equation (8.9)]{AFK} to the three representative changes and
multiply by the phase in AFK Definition~1.30.  Direct substitution gives
\[
 \Phi_{p,q}\,(\hbox{three translation multipliers})
 =(-1)^{E(p,q)},\qquad
 E(p,q)=5(p+q)+pq+\min(5,p+q).
\]
This is a finite parity calculation:
\[
 E(p,q)\equiv
 \begin{cases}
 pq& (p+q\le5),\\
 p+q+pq+1&(p+q\ge5)
 \end{cases}\pmod2,
\]
where the two expressions agree at $p+q=5$.  In the order
$p=0,\ldots,4$ (columns $q=0,\ldots,4$), the resulting signs are
\[
\begin{pmatrix}
 *&+&+&+&+\\
 +&-&+&-&+\\
 +&+&+&+&-\\
 +&-&+&+&+\\
 +&+&-&+&-
\end{pmatrix};
\]
the asterisk is the excluded zero characteristic.  This accounts for
all twenty-four cases and proves (2) directly from the cited
definitions.
\end{proof}

\section{Ray labels and all characteristics}

\begin{lemma}[Certified ray arithmetic]\label{lem:ray}
The class and ray groups are
\[
 \Cl(K)=1,\qquad
 \Cl_{(5)}(K)\cong C_4,\qquad
 \Cl_{\mathfrak m}(K)\cong C_8,\qquad
 \Cl_{(5)\infty_1\infty_2}(K)\cong C_8\times C_2.
\tag{3}
\]
Kopp's residue/sign class $\mathfrak R$ is the element $4\in C_8$.
The fiber of
$\Cl_{(5)\infty_1\infty_2}(K)\to\Cl_{\mathfrak m}(K)$ has order two,
so the exponent in Kopp's theorem is $n=2/2=1$.
The computation is unconditional.
\end{lemma}

\begin{proof}
PARI/GP~2.15.4 gives the structures in (3), and
\texttt{bnfcertify}=1.  The unit group is
$\{\pm\beta^n:n\in\Z\}$ and $\beta^3\equiv1\pmod5$.
Hence units congruent to $-1$ modulo $5$ are
$-\beta^{3k}$.  Their value under $\infty_2$ is
$-(2-\sqrt3)^{3k}<0$, so no global unit realizes the pair
$(-1\bmod5,+\text{ at }\infty_2)$ defining Kopp's $\mathfrak R$.
It is therefore the nontrivial sign class, namely the unique element
of order two in $C_8$.
\end{proof}

\begin{proposition}[No quadratic-character shortcut]\label{prop:noquadratic}
The Kopp partial-zeta difference has zero coefficient at the unique
order-two character of $\Cl_{\mathfrak m}(K)$.  It is supported on four
characters, all of order eight.
\end{proposition}

\begin{proof}
Write $\Cl_{\mathfrak m}(K)=\langle g\rangle\cong C_8$ and
$\mathfrak R=g^4$.  For
$\chi_k(g)=\exp(2\pi i k/8)$, Fourier inversion gives
\[
\zeta_{\mathfrak m}(s,\mathfrak A)
-\zeta_{\mathfrak m}(s,\mathfrak R\mathfrak A)
=\frac18\sum_{k=0}^7
\overline{\chi_k(\mathfrak A)}
\bigl(1-\overline{\chi_k(\mathfrak R)}\bigr)L(s,\chi_k).
\tag{CF}
\]
Since $\chi_k(\mathfrak R)=(-1)^k$, the coefficient is zero for even
$k$ and two for odd $k$.  The support is therefore
$k=1,3,5,7$, whose characters all have order eight.  The unique
quadratic character is $\chi_4$; it is even on $\mathfrak R$ and its
coefficient in \textup{(CF)} is zero.  Hence the difference cannot
factor through the quadratic quotient.  This is the precise reason
the dimension-four quadratic $L$-function reduction does not recur.
Equivalently, $\mathfrak R=g^4$ lies in the kernel
$\langle g^2\rangle$ of the quotient $C_8\to C_2$, so
$\mathfrak A$ and $\mathfrak R\mathfrak A$ have the same image in the
unique quadratic quotient and their difference projects to zero.
\end{proof}

For each $(p,q)\ne(0,0)$ start with $p$ and repeatedly subtract five;
let $\widetilde p$ be the first lift for which
\[
 q(2-\sqrt3)-\widetilde p>0.
\tag{4}
\]
Let $\mathfrak q$ be the prime above $3$, and use it as the generator
of $C_8$.  PARI's raw cyclic basis $g$ gives the discrete logarithms
of the positive principal ideals $(q\beta-\widetilde p)$ in the audit
transcript, with $\mathfrak q=g^5$.  Multiplication by $5$, which is
its own inverse modulo $8$, converts those raw logs to the normalized
$\mathfrak q$-coordinates below:
\[
\begin{array}{c|ccccc}
p\backslash q&0&1&2&3&4\\ \hline
0&-&0&6&2&4\\
1&4&7&5&7&0\\
2&2&5&5&6&3\\
3&6&7&2&1&1\\
4&0&4&3&1&3
\end{array}
\tag{5}
\]
The corresponding positive lifts are
$0$ in the first row, $-4$ in row one except at $(1,4)$ where it is
$1$, and $-3,-2,-1$ in rows two, three, and four.

\begin{lemma}[All multiplier comparisons]\label{lem:multiplier}
For every nonzero $(p,q)\bmod5$, Kopp's modular multiplier equals
the square of the AFK phase at that characteristic.
\end{lemma}

\begin{proof}
Write $B=(\begin{smallmatrix}a&b\\c&d\end{smallmatrix})$ and
$\boldsymbol r=(\widetilde p/5,q/5)^{\mathsf T}$, with
$\widetilde p$ chosen by (4).  Substitution in Kopp's defining
theta-character gives the exponent
\[
\frac12\{(c-d+1)r_1+(-a+b+1)r_2-cd r_1^2
+2(a-1)d r_1r_2-(a-2)b r_2^2\}.
\tag{6}
\]
For $B=(\begin{smallmatrix}56&-15\\15&-4\end{smallmatrix})$, substituting
the particular positive lift from (4) and reducing (6) modulo one gives
\[
 \frac{p^2-4pq+q^2}{5}\pmod1
\tag{7}
\]
separately for each of the twenty-four nonzero residue pairs $(p,q)$.
Thus (7) is a certified finite case calculation with the stated lifts,
not a polynomial identity for arbitrary integral $p,q$.  Also
$s(-4,15)=-19/90$, so the Rademacher invariant is three and
$\psi^2(B)=i$.  Therefore the Kopp multiplier has exponent
\[
 -\frac14-\frac{p^2-4pq+q^2}{5}\pmod1,
\tag{8}
\]
which is precisely the exponent of the square of the AFK phase.
The exact twenty-four-row substitution, including every positive lift,
is part of the bridge certificate.
\end{proof}

For a normalized ray log $k$, name the positive square root of its
ray value by
\[
\begin{array}{c|cccccccc}
k&0&1&2&3&4&5&6&7\\ \hline
\text{positive root}&x&w&y^{-1}&z&x^{-1}&w^{-1}&y&z^{-1}.
\end{array}
\]
The class log in (5) selects the unsigned entry and the explicit sign
in (2) selects its sign.  Thus the complete reconstruction table is
\[
T_5=
\begin{pmatrix}
\sqrt6&x&y&y^{-1}&x^{-1}\\
x^{-1}&-z^{-1}&w^{-1}&-z^{-1}&x\\
y^{-1}&w^{-1}&w^{-1}&y&-z\\
y&-z^{-1}&y^{-1}&w&w\\
x&x^{-1}&-z&w&-z
\end{pmatrix}.
\tag{9}
\]
Here $x,y,z,w$ denote the four positive representatives of
the ray-class pairs
\[
(0,4),\qquad(6,2),\qquad(3,7),\qquad(1,5),
\tag{10}
\]
respectively.  In particular, the table is generated from the ray logs
and the sign exponent; it is not stored as a list of expected answers in
the bridge generator.

\section{The analytic ray packet}

Set
\[
\begin{aligned}
P(U)={}&U^8-(8+5\sqrt3)U^7+(53+30\sqrt3)U^6\\
&-(156+90\sqrt3)U^5+(225+130\sqrt3)U^4\\
&-(156+90\sqrt3)U^3+(53+30\sqrt3)U^2\\
&-(8+5\sqrt3)U+1.
\end{aligned}
\tag{11}
\]
It is the reciprocal lift $P(U)=U^4F(U+U^{-1})$ of
\[
\begin{aligned}
F(V)={}&V^4-(8+5\sqrt3)V^3+(49+30\sqrt3)V^2\\
&-(132+75\sqrt3)V+(121+70\sqrt3).
\end{aligned}
\tag{12}
\]
Kopp previously recognized the same polynomial numerically as the
Stark-unit polynomial for $\Q(\sqrt3)$ and modulus five; see
\cite[\S7, (7.21)]{Kopp2021} and
\cite[Example~1.17, (3.12)]{Kopp2023}.  PARI's \texttt{bnrstark}
reproduces (12), but that recognition is not used as a proof below.

The square roots of the reciprocal packet have absolute polynomial
\[
\begin{aligned}
\mathcal Q(X)={}&X^{32}-16X^{30}+95X^{28}-260X^{26}+355X^{24}
-348X^{22}\\
&+388X^{20}-300X^{18}+195X^{16}-300X^{14}+388X^{12}\\
&-348X^{10}+355X^8-260X^6+95X^4-16X^2+1.
\end{aligned}
\tag{13}
\]
Define $w^{-1},z^{-1},x^{-1},y^{-1},y,x,z,w$ to be the unique
positive roots of (13) in, respectively,
\[
\begin{gathered}
(0.424835,0.424836),\ (0.481476,0.481477),\
(0.506963,0.506964),\ (0.782390,0.782391),\\
(1.278133,1.278134),\ (1.972526,1.972527),\
(2.076946,2.076947),\ (2.353849,2.353850).
\end{gathered}
\tag{14}
\]
Sturm's theorem certifies exactly one root in every interval.

\subsection{Shintani's theorem with an explicit exponent}

For a ray class $A\in\Cl_{\mathfrak m}(K)$ put
\[
 X_A:=\exp Z'_{\mathfrak m}(0,A).
\]

\begin{lemma}[Specialized Shintani algebraicity]\label{lem:shintani}
For every $A\in\Cl_{\mathfrak m}(K)$,
\[
 X_A^{5760}\in\mathcal O_H^\times,
\]
where $H$ is the ray field of modulus $\mathfrak m$.  These powers have
the Artin action indexed by the ray classes.  At every embedding of
$H$ inducing $\infty_2$ on $K$, their absolute value is one.
\end{lemma}

\begin{proof}
Work first in the full narrow ray group
\[
 C:=\Cl_{(5)\infty_1\infty_2}(K)\cong C_8\times C_2.
\]
In Shintani's notation, let $\mu$ be the class represented by
$11-5\beta$, which is congruent to $1$ modulo $5$ and has signs
$(-,+)$ at $(\infty_1,\infty_2)$.  Let $\nu$ be the class represented
by $4$, which is congruent to $-1$ modulo $5$ and is totally positive.
The exact ray logs are
\[
 \mu=[4,1]\in C,\qquad \nu=[4,0]\in C.
\]
Forgetting $\infty_1$ kills $\mu$ and sends $\nu$ to the unique
order-two element $4\in C_8$.  Hence
\[
 G=\langle\mu\rangle,\qquad \nu\notin G,\qquad
 C/G\simeq\Cl_{\mathfrak m}(K).
\]

Every unit is $\pm\beta^n$, and $\beta$ has order three modulo $5$.
Thus no totally positive unit is $-1$ modulo $5$, which is Shintani's
condition (0-3).  There is no norm-$-1$ unit, since
$a^2-3b^2=-1$ is impossible modulo $3$; this implies condition (0-6).

The field $H$ has signature $(8,4)$ and a unique octic subfield
\[
 M\simeq\Q(\zeta_{60})^+.
\]
An exhaustive degree-eight subfield enumeration of the degree-sixteen
defining field returns exactly one subfield; an exact
\texttt{nfisisom} calculation then identifies it with the displayed
cyclotomic field.  This is recorded in the archived
\path{scripts/dimension_five_shintani_audit.gp} transcript and checked
independently in
\path{scripts/dimension_five_special_reduction_audit.gp}.
Thus $H/M$ is quadratic, $M$ is the maximal absolutely abelian
subfield of $H$, and exactly one real place of $K$ splits in $H$.
These are Shintani's hypothesis (0-9).  His Theorem~2
\cite{Shintani1978} therefore applies to $X_{(5)}(c,G)$.

This invariant is exactly the one-place zeta difference used here.
Indeed, if $c$ maps to $A$, then the inverse images of $A$ and
$RA$ in $C$ are respectively
\[
 \{c,c\mu\},\qquad \{c\nu,c\mu\nu\}.
\]
Shintani's Theorem~1 consequently gives
\[
\begin{aligned}
\log X_{(5)}(c,G)
&=\zeta'_{\rm full}(0,c)-\zeta'_{\rm full}(0,c\nu)\\
&\quad+\zeta'_{\rm full}(0,c\mu)
       -\zeta'_{\rm full}(0,c\mu\nu)\\
&=\zeta'_{\mathfrak m}(0,A)-\zeta'_{\mathfrak m}(0,RA)
=Z'_{\mathfrak m}(0,A).
\end{aligned}
\]

It remains to identify the imaginary quadratic field in Shintani's
proof and to clear every exponent in his construction.  Write
$\mathcal N$ for the full two-infinite-place ray field, the normal
closure of $H$.  In the cyclic coordinates
\[
 C=\langle e_1\rangle\times\langle e_2\rangle
 \cong C_8\times C_2
\]
returned by the exact ray computation, the nontrivial automorphism
$\sqrt3\mapsto-\sqrt3$ acts by
\[
 e_1\longmapsto(5,0),\qquad e_2\longmapsto(4,1).
\tag{S1}
\]
Consequently its fixed subgroup is
\[
 C^{\mathrm{fix}}
 =\{(a,b):a+b\equiv0\pmod2\}
 =\langle(2,0),(1,1)\rangle.
\tag{S2}
\]
Here $G_1$ in Shintani's proof is trivial.  His index-two subgroup
$(C/G_1)_0$ has the property that $c^\sigma c^{-1}\in G_1$ for each
of its classes, so it is contained in (S2); since both have index two,
they are equal.  Exact class-field
construction with the HNF matrix
$\left(\begin{smallmatrix}2&1\\0&1\end{smallmatrix}\right)$ gives
\[
 \mathcal N^{C^{\mathrm{fix}}}
 =K(\sqrt{-5}).
\tag{S3}
\]
Shintani's field $L$ is therefore the biquadratic field
$K\,k$ with
\[
 k=\Q(\sqrt{-5}).
\]
This derives the imaginary quadratic field occurring in his proof;
it does not select one by numerical recognition from a list of
quadratic subfields.

For an independent class-field check, the ray group of $k$ modulo
$15$ is $C_{40}\times C_2$.  Killing the subgroup with HNF matrix
$\operatorname{diag}(8,2)$ produces a degree-sixteen extension of
$k$.  Its exact conductor, in the integral basis used by PARI, is
\[
 \mathfrak c=
 \begin{pmatrix}15&0\\0&3\end{pmatrix},
\tag{S4}
\]
and \texttt{nfisisom} identifies its absolute defining field with
$\mathcal N$.  Thus it is the ray-class subfield over the very same
$k$ forced by (S1)--(S3).  The computation also gives
$h_k=2$ and \texttt{bnfcertify}=1.

We now spell out the exponent bookkeeping.  Factor
\[
 \mathfrak c=\mathfrak p_3\overline{\mathfrak p}_3\mathfrak p_5,
\quad
 \mathfrak p_3=\begin{pmatrix}3&2\\0&1\end{pmatrix},\quad
 \overline{\mathfrak p}_3=\begin{pmatrix}3&1\\0&1\end{pmatrix},\quad
 \mathfrak p_5=\begin{pmatrix}5&0\\0&1\end{pmatrix}.
\]
For $\mathfrak d\mid\mathfrak c$, let $f_{\mathfrak d}$ be the
smallest positive rational integer in $\mathfrak d$, let
$h(\mathfrak d)$ be the order of its ray group, and let
$w(\mathfrak d)$ be the number of roots of unity in $k$ congruent to
$1$ modulo $\mathfrak d$.  Here $h(\mathfrak c)=16$.  Put, exactly as
in Shintani's equation~(8),
\[
 n(\mathfrak d)
 =w(\mathfrak d)\frac{h(\mathfrak c)}{h(\mathfrak d)}.
\tag{S5}
\]
The exact divisor audit is
\[
\begin{array}{c@{\quad}c@{\quad}c@{\quad}c@{\quad}c@{\quad}c}
\mathfrak d&f_{\mathfrak d}&h(\mathfrak d)&w(\mathfrak d)
 &n(\mathfrak d)&\text{clearing exponent}\\ \hline
\mathcal O_k&1&2&2&16&12h_kn(\mathfrak d)=384\\
\mathfrak p_3&3&2&1&8&12f_{\mathfrak d}n(\mathfrak d)=288\\
\overline{\mathfrak p}_3&3&2&1&8&288\\
(3)&3&4&1&4&144\\
\mathfrak p_5&5&4&1&4&240\\
\mathfrak p_3\mathfrak p_5&15&8&1&2&360\\
\overline{\mathfrak p}_3\mathfrak p_5&15&8&1&2&360\\
\mathfrak c&15&16&1&1&180
\end{array}
\tag{S6}
\]
Indeed, Shintani's equations~(4) and (9) write, for
$\mathfrak d\ne\mathcal O_k$,
\[
 Z_{\mathfrak d}=|\varphi_{\mathfrak d}|^{1/(6f_{\mathfrak d})},
 \qquad
 W_{\mathfrak c}
 =\prod_{\mathfrak d\mid\mathfrak c}
  |Z_{\mathfrak d}(\cdots)|^{1/n(\mathfrak d)}.
\]
Replacing an absolute value by the product with its complex conjugate
introduces the additional factor two, so
$12f_{\mathfrak d}n(\mathfrak d)$ clears every denominator.  For
$\mathfrak d=\mathcal O_k$, Shintani's equation~(6) has exponent
$1/(6h_k)$ instead, giving $12h_kn(\mathfrak d)$.  Proposition~4
is, literally in his notation,
\[
 Y_{(5)}(c,G)
 =\prod_{h\in H_1}
   \frac{W_{\mathfrak c}(\widetilde c\,\nu h)}
        {W_{\mathfrak c}(\widetilde c\,h)}.
\tag{S7}
\]
Thus every $W$-invariant occurs with integral exponent $+1$ or $-1$;
there is no additional rational exponent.  The least common multiple
of the last column of
(S6) is
\[
 \operatorname{lcm}(384,288,288,144,240,360,360,180)=5760.
\tag{S8}
\]
Finally, $(5)$ is prime and its only proper divisor fails Shintani's
condition (0-3), so his $Y_{(5)}$ equals $X_{(5)}$ and no induction
exponent is added.  Proposition~5 now supplies the unit, Artin, and
nonsplit-modulus-one assertions.
\end{proof}

\subsection{The labeled candidate field}

The absolute polynomial of a root of (11) is
\[
\begin{aligned}
p(T)={}&T^{16}-16T^{15}+95T^{14}-260T^{13}+355T^{12}
-348T^{11}\\
&+388T^{10}-300T^9+195T^8-300T^7+388T^6-348T^5\\
&+355T^4-260T^3+95T^2-16T+1.
\end{aligned}
\]

\begin{lemma}[Exact field and Frobenius labels]\label{lem:frobenius}
The field defined by (11) is the labeled ray field $H$ over $K$.
Let $U$ be the root of $p$ in
\[
 (3.890,3.891).
\]
Then $U$ is a norm-one unit.  Arithmetic Frobenius at the prime
$\mathfrak q$ above $3$ acts transitively on its eight real conjugates,
which lie successively in
\[
\begin{gathered}
(3.890,3.891),\ (5.540,5.541),\ (0.612,0.613),\
(4.313,4.314),\\
(0.257,0.258),\ (0.180,0.181),\ (1.633,1.634),\
(0.231,0.232).
\end{gathered}
\]
At the embeddings above $\infty_2$, all conjugates of $U$ have
absolute value one.
\end{lemma}

\begin{proof}
Applying \texttt{rnfequation} to both the candidate and the ray field,
\texttt{nfisisom} returns eight maps.  Exact substitution shows that
all eight maps fix the labeled root of $y^2-4y+1$; hence the
isomorphism is over $K$, not merely over $\Q$.  Moreover
\texttt{bnfcertify}=1, \texttt{bnfisunit} gives an exact unit
decomposition, and the norm of $U$ is one.

The prime $\mathfrak q$ generates the ray group.  Modulo
$\mathfrak q$ the relative polynomial is
\[
 T^8+T^7+2T^6+2T^2+T+1\in\mathbb F_3[T],
\]
which is irreducible.  Arithmetic Frobenius is therefore $T\mapsto
T^3$.  Reduction modulo this polynomial identifies one of the exact
\texttt{nfgaloisconj} automorphisms with Frobenius.

The certificate refines $(3.890,3.891)$ to a rational subinterval of
width $10^{-80}$ containing exactly one root of $p$.  Exact rational
interval evaluation of the
successive Frobenius polynomials places their images inside the eight
disjoint intervals in the statement.  This proves the labels
Galois-equivariantly.

Finally, under the nonidentity embedding of $K$, substitute
$T=2(Y^2-1)/(Y^2+1)$ in the trace polynomial (12).
The exact \texttt{nfpolsturm} counts are $[8,0]$, so all four trace
roots at that embedding lie in $(-2,2)$.  Each corresponding
reciprocal quadratic has both roots on the unit circle.
\end{proof}

\subsection{Certified removal of the power}

\begin{theorem}[Unconditional ray packet]\label{thm:packet}
For
\[
 \mathcal X_k:=\exp Z'_{\mathfrak m}(0,\mathfrak q^k),
 \qquad 0\le k<8,
\]
the exact Artin-labeled packet is
\[
(\mathcal X_0,\ldots,\mathcal X_7)
=(x^2,w^2,y^{-2},z^2,x^{-2},w^{-2},y^2,z^{-2}).
\tag{15}
\]
In particular, $\mathcal X_0=U$.
\end{theorem}

\begin{proof}
Kopp's proved Kronecker-limit theorem, with the convention and
multiplier checked above, gives
\[
 \widetilde\nu_{p,q}^{\,2}
 =\exp Z'_{\mathfrak m}(0,\mathfrak A_{p,q});
\tag{16}
\]
here $\mathfrak A_{p,q}$ is the normalized ray class in (5).
The four positive generators at characteristics
\[
 (0,1),\quad(0,2),\quad(2,4),\quad(3,3)
\]
have $\mathfrak q$-coordinates $0,6,3,1$, respectively.

We evaluate the regularized integral for the reciprocal double sine
(1) with Arb balls through \texttt{python-flint}~\cite{pythonflint}.
Every Simpson panel uses an Arb enclosure of the
fourth derivative.  All endpoints, widths, tolerances, and acceptance
comparisons are exact rationals or outward-rounded balls.  Near zero,
the positive series of
$H_c(t)=\sinh(ct/2)/(ct/2)$ gives the rational bound
\[
 |f(t)-f(0)|\le4t^2\qquad(0\le t\le10^{-4}).
\]
For the tail, every reduced argument and its complement belongs to
$(1/100,5)$; beyond $36$, each exponential denominator exceeds
$3/4$, giving the remainder bound $e^{-36}/9$.
The candidate roots are selected by disjoint rational isolating
intervals, independently of library ordering.

For each of the four characteristics, the Arb ball for
$\log(\widetilde\nu_{p,q}^{\,2})-\log(U_k)$ contains zero.  The largest
upper bound for the absolute value of any of these four balls is
\[
 \epsilon=4.4\times10^{-11}.
\]
These four enclosures are logically needed before the packet has been
identified: they cover one representative of each reciprocal pair of
real ray classes, and their reciprocal partners bound the other four
real embeddings.  Once equality at one class is proved, the exact Artin
action propagates it through the packet.  Thus the extra three
enclosures certify the eight-place height bound and the characteristic
labels; they do not replace any step of Galois propagation.

By Lemma~\ref{lem:shintani}, both
$\mathcal X_0^{5760}$ and $U^{5760}$ are units in the same labeled
field.  Set
\[
 \eta=\mathcal X_0^{5760}/U^{5760}.
\]
Put
\[
 \delta:=5760\epsilon<2.54\times10^{-7}.
\]
Lemma~\ref{lem:frobenius} and the four reciprocal-pair enclosures show
that at every one of the eight real embeddings,
\[
 |\log|\sigma(\eta)||<\delta.
\]
At the four complex pairs, both numerator and denominator have
absolute value one.  Since $H$ has signature $(8,4)$ and degree
sixteen, and since the finite-place terms vanish because $\eta$ is a
unit, the absolute logarithmic Weil height is therefore
\[
\begin{aligned}
 h(\eta)
 &=\frac1{16}\left(
 \sum_{\sigma\ {\rm real}}\max(0,\log|\sigma\eta|)
 +2\sum_{\tau\ {\rm complex\ pair}}\max(0,\log|\tau\eta|)
 \right)\\
 &<\frac8{16}\delta
 <1.27\times10^{-7}.
\end{aligned}
\tag{H}
\]

If $3\le d=[\Q(\eta):\Q]\le16$ and $\eta$ is not a root of unity,
Voutier's explicit bound \cite{Voutier1996} gives
\[
 h(\eta)>
 \frac1{4d}\left(\frac{\log\log d}{\log d}\right)^3.
\]
Its minimum for $3\le d\le16$ is
$5.2279533222\ldots\times10^{-5}$, a contradiction.  A non-torsion
quadratic unit has height at least
$\tfrac12\log((1+\sqrt5)/2)$, and a rational unit is $\pm1$.
Thus $\eta$ is a root of unity.  Since $H$ has a real embedding, its
only roots of unity are $\pm1$; positivity at the distinguished
embedding gives $\eta=1$.  Therefore
$\mathcal X_0^{5760}=U^{5760}$ and positivity yields
$\mathcal X_0=U$.  The exact Artin action in
Lemma~\ref{lem:shintani}, together with the certified Frobenius orbit
in Lemma~\ref{lem:frobenius}, proves the full packet (15).
\end{proof}

Since every reciprocal double-sine factor in (2) is positive at the
displayed arguments, equations (2), (5), and (15) now derive the
signed table (9) unconditionally.

\section{Exact finite cancellation}

Let $\zeta_5=e^{2\pi i/5}$ and
$\tau=-e^{\pi i/5}=\zeta_5^3$.  Define, without reference to software,
\[
 \Pi^{(5)}_{r,c}:=\frac1{5\sqrt6}
 \sum_{\substack{0\le a,b<5\\r\equiv c+a\pmod5}}
 T_{a,b}\tau^{ab}\zeta_5^{bc},
 \qquad 0\le r,c<5,
\tag{17}
\]
where $T$ is (9).  Equivalently,
$\Pi^{(5)}=\frac15\sum_{a,b}(T_{a,b}/\sqrt6)D_{a,b}$ for
$D_{a,b}e_c=\tau^{ab}\zeta_5^{bc}e_{c+a}$.
In particular,
\[
 \Tr\Pi^{(5)}=T_{0,0}/\sqrt6=1.
\tag{18}
\]

\begin{proposition}[Exact 100-minor certificate]\label{prop:minors}
For the algebraic values selected by (13)--(14), all one hundred
$2$-minors of the matrix (17) vanish exactly.  Hence
$\operatorname{rank}\Pi^{(5)}=1$ and
$(\Pi^{(5)})^2=\Pi^{(5)}$.
\end{proposition}

\begin{proof}
Over the coefficient field $\Q(\zeta_5,\sqrt6)$, (13) has four
degree-eight factors.  Let $w$ be the interval-defined root in (14).
Exact rational interval arithmetic applied to the
\texttt{nfgaloisconj} polynomials proves that
\texttt{nfgaloisconj} substitutions are
\[
\begin{array}{c|cccccccc}
 &w&w^{-1}&y&y^{-1}&x&x^{-1}&z&z^{-1}\\ \hline
\text{conjugate index}&2&8&5&4&10&6&16&12.
\end{array}
\tag{19}
\]
with every image lying inside its corresponding interval in (14).
Moreover, exact \texttt{nfisincl} identities express $\sqrt5$ and
$\sqrt6$ as polynomials in $w$; rational interval evaluation selects
their positive signs.  Among the four factors, factor four is the
unique factor on which both of these subfield elements equal the
positive $\sqrt5$ and positive $\sqrt6$.  This certifies, rather than
numerically assumes, that the interval-defined tuple is the tuple used
in the factor-four calculation.

Substitution in (9) and (17), followed by exact reduction modulo the
selected factor, makes all one hundred minors zero.  As a factor
selection check, factors one, two, and three leave respectively
$70,100,100$ nonzero minors.  Factor one preserves the positive
$\sqrt5$ embedding but reverses $\sqrt6$; thirty minors vanish for
structural reasons common to those two embeddings, while the remaining
seventy detect the wrong $\sqrt6$ sign.  Factors two and three already
reverse $\sqrt5$ and leave all one hundred minors nonzero.  Thus the
factor check is not circular.  This is exact number-field arithmetic,
not floating-point evaluation.  The complete script and transcript
are supplied with the manuscript.

The matrix is nonzero because (18) gives trace one.  Thus its rank is
exactly one.  Any rank-one matrix $M$ satisfies
$M^2=(\Tr M)M$, proving idempotency.
\end{proof}

\section{Dimension-five twisted convolution}

\begin{lemma}[Weyl coefficient bridge]\label{lem:bridge}
The matrix $\Pi^{(5)}$ of (17) is the standard-basis Weyl reconstruction of
the canonical AFK ghost for
$t=(5,1,\langle1,-4,1\rangle)$.  Moreover,
$(\Pi^{(5)})^2=\Pi^{(5)}$ is equivalent to
the $\lambda=1$ instance of AFK equation~(1.49):
\[
\sum_{\boldsymbol q\in\mathcal I_{\boldsymbol p}}
\zeta_5^{\langle\boldsymbol p,(I+L_{z,t})\boldsymbol q\rangle}
\operatorname{shin}_{A_t}^{\boldsymbol q/5}(\beta)
\operatorname{shin}_{A_t^{-1}}^{(\boldsymbol q-\boldsymbol p)/5}(\beta)
=25\delta_{\boldsymbol p,\boldsymbol0}^{(5)}.
\tag{20}
\]
\end{lemma}

\begin{proof}
For $\boldsymbol p\ne\boldsymbol0$, put
$\nu_{\boldsymbol p}:=\widetilde\nu_{\boldsymbol p}$ as given by
(2).  The first two unconditional parts of AFK's construction give
\[
 \nu_{\boldsymbol p}\nu_{-\boldsymbol p}=1.
\tag{21}
\]
Theorem~\ref{thm:packet}, followed by (2) and (5), identifies these
twenty-four numbers with the nonzero entries of $T$.  The actual ghost
reconstruction is therefore
\[
 \Pi^{(5)}=\frac15I+\frac1{5\sqrt6}
 \sum_{\boldsymbol p\ne\boldsymbol0}
 \nu_{\boldsymbol p}D_{\boldsymbol p}.
\tag{22}
\]
The artificial entry $T_{0,0}=\sqrt6$ packages the first term of
(22) into the single entry formula (17); it is not an analytic
zero-characteristic value.

Now expand $(\Pi^{(5)})^2-\Pi^{(5)}$ using
$D_{\boldsymbol p}D_{\boldsymbol q}
=\tau^{\langle\boldsymbol p,\boldsymbol q\rangle}
D_{\boldsymbol p+\boldsymbol q}$ and compare every Weyl coefficient.
For $\boldsymbol p\ne\boldsymbol0$, multiplication by $150$ makes its
coefficient
\[
 \sum_{\substack{\boldsymbol q\bmod5\\
 \boldsymbol q\ne\boldsymbol0,\boldsymbol p}}
 \tau^{\langle\boldsymbol p,\boldsymbol q\rangle}
 \nu_{\boldsymbol p-\boldsymbol q}\nu_{\boldsymbol q}
-3\sqrt6\,\nu_{\boldsymbol p}.
\tag{23}
\]
Using (21), the product in (23) is
$\nu_{\boldsymbol q}/\nu_{\boldsymbol q-\boldsymbol p}$.

AFK's $\nu_{\boldsymbol0}^{\rm aux}$ is an auxiliary cocycle value,
not the identity coefficient in (22).  Their zero-characteristic
lemma specializes to
\[
 \nu_{\boldsymbol0}^{\rm aux}
+(\nu_{\boldsymbol0}^{\rm aux})^{-1}
=-(5-2)\sqrt6=-3\sqrt6.
\tag{24}
\]
After adjoining the two omitted endpoints
$\boldsymbol q=\boldsymbol0,\boldsymbol p$, their combined contribution
is
\[
 \nu_{\boldsymbol p}\!
 \left(\nu_{\boldsymbol0}^{\rm aux}
+(\nu_{\boldsymbol0}^{\rm aux})^{-1}\right)
=-3\sqrt6\,\nu_{\boldsymbol p}.
\]
Thus (23) is exactly the full endpoint-corrected sum
\[
 \sum_{\boldsymbol q\in\mathcal I_{\boldsymbol p}}
 \tau^{\langle\boldsymbol p,\boldsymbol q\rangle}
 \frac{\nu_{\boldsymbol q}}
 {\nu_{\boldsymbol q-\boldsymbol p}},
\tag{25}
\]
with the two endpoint ratios interpreted using
$\nu_{\boldsymbol0}^{\rm aux}$.

Write $\Phi_{\boldsymbol q}$ for the AFK phase.  Their phase-ratio
calculation specializes, with $r=f_j/f=1$ and $d_j=d=5$, to
\[
 \frac{\Phi_{\boldsymbol q}}
 {\Phi_{\boldsymbol q-\boldsymbol p}}
 =
 (-1)^{1+s_5(\boldsymbol p)}
 \tau^{Q(\boldsymbol p)}
 \tau^{\langle\boldsymbol p,\boldsymbol q\rangle}
 \zeta_5^{\langle\boldsymbol p,L_{z,t}\boldsymbol q\rangle}.
\tag{26}
\]
The first two factors are nonzero and independent of
$\boldsymbol q$.  Multiplying the remaining $\tau$ factor by the
displacement chirp in (25), and using $\tau^2=\zeta_5$, gives exactly
\[
 \zeta_5^{\langle\boldsymbol p,
 (I+L_{z,t})\boldsymbol q\rangle}.
\]
The cocycle inverse identity
\[
 \frac{\operatorname{shin}_{A_t}^{\boldsymbol q/5}(\beta)}
 {\operatorname{shin}_{A_t}^{(\boldsymbol q-\boldsymbol p)/5}(\beta)}
 =
 \operatorname{shin}_{A_t}^{\boldsymbol q/5}(\beta)
 \operatorname{shin}_{A_t^{-1}}^{
   (\boldsymbol q-\boldsymbol p)/5}(\beta)
\tag{27}
\]
gives the two factors printed in (20).  Here
$A_t=B$, $L_{z,t}=L$, and
$f_t(1)=2+9\equiv1\pmod5$, so the identity twist produces exactly the
chirp in (25).  The remaining factor depends on
$\boldsymbol p$ but is nonzero and independent of
$\boldsymbol q$.  Hence the vanishing of every nonzero Weyl coefficient
is equivalent to (20) for $\boldsymbol p\ne\boldsymbol0$.
For $\boldsymbol p=\boldsymbol0$, the cocycle inverse identity makes
the sum $25$ automatically.  This is the specialization of AFK
equations (5.78)--(5.98), including rather than suppressing their
endpoint correction.
\end{proof}

\begin{theorem}[Complete dimension-five TCC]\label{thm:tcc}
For every dimension-five admissible tuple $t$,
\[
 0,1\in\mathcal Z_t.
\]
Furthermore $\mathcal Z_{t'}=\mathcal Z_t$ for any two dimension-five
admissible tuples $t,t'$.
\end{theorem}

\begin{proof}
Proposition~\ref{prop:minors} and Lemma~\ref{lem:bridge} prove the
$\lambda=1$ identity for the canonical tuple.  The twist condition is
\[
 \det(G)f_t(\lambda)\equiv1\pmod5,\qquad
 f_t(\lambda)=2\lambda+9.
\]
Thus $G=I$ gives $2\lambda+9\equiv1\pmod5$, or
$\lambda\equiv1\pmod5$.

Complex conjugation preserves idempotency and sends
$\lambda$ to
$\bar\lambda=-(\lambda+d-1)=1-\lambda\pmod5$.
It therefore sends the verified shift $1$ to shift $0$.
Equivalently the conjugate twist has determinant $-1$, and
$-(2\lambda+9)\equiv1\pmod5$ gives $\lambda\equiv0\pmod5$.

Finally, dimension-five admissibility forces $r=1$,
$K=\Q(\sqrt3)$, and $j=m=1$.  Since the relevant conductor is one,
every admissible form has discriminant twelve.  The ordinary class
number is one, so primitive indefinite forms of discriminant twelve
constitute one $\mathrm{GL}_2(\Z)$-equivalence class.  (The narrow class
number is not needed: determinant-$-1$ transformations are allowed.)
Thus every admissible form is integrally equivalent to the canonical
one in the sense required here.

For clarity, the transport does not invoke AFK's conditional
transformation theorem.  If $t_M$ is obtained from $t$ by
$M\in\mathrm{GL}_2(\Z)$, then unconditionally
\[
 A_{t_M}=M^{-1}A_tM,\qquad
 L_{z,t_M}=M^{-1}L_{z,t}M,\qquad
 \beta_{t_M}=M^{-1}\beta_t.
\]
Kopp's covariance formula sends each cocycle factor at
$(t_M,\boldsymbol q)$ to the corresponding factor at
$(t,\ell M\boldsymbol q)$, with
$\ell=\operatorname{sgn}(j_{M^{-1}}(\beta_t))$; when
$\det M=-1$ both factors are complex-conjugated, which does not change
the real TCC sum.  Reindexing
\[
 \boldsymbol q'=\ell M\boldsymbol q,\qquad
 \mathcal I'_{\ell M\boldsymbol p}=\ell M\mathcal I_{\boldsymbol p},
\]
and using symplectic covariance transforms the complete left side of
(20) for $(t_M,\boldsymbol p)$ into that for
$(t,\ell M\boldsymbol p)$.  Hence
$\mathcal Z_t\subseteq\mathcal Z_{t_M}$.  Applying the same calculation to
$M^{-1}$ gives the reverse inclusion and therefore equality.
\end{proof}

\begin{proof}[Proof of Theorem~\ref{thm:lowdim}]
The dimension-four assertion is proved in Section~\ref{sec:d4}; the
dimension-five assertion is
Theorem~\ref{thm:tcc}.
\end{proof}

\section{Composite dimensions six and eight}\label{sec:composite}

The preceding argument succeeds because Shintani's 1978 theorem supplies
powered algebraicity for the entire dimension-five packet.  The same
finite reconstruction can be completed in dimensions six and eight, but
none of the special-value theorems used in this paper supplies the
needed oriented algebraicity for every primitive characteristic.  This
section records computationally exact finite results, the unconditional
analytic strata covered by the cited theorems, and the remaining
research targets.  It contains no dimension-six or dimension-eight TCC
theorem.

\subsection{Dimension six: one primitive order-six value}

Put
\[
 K_6=\Q(\sqrt{21}),\qquad
 \beta_6=\frac{5+\sqrt{21}}2,\qquad
 Q_6=\langle1,-5,1\rangle.
\]
The relevant one-place ray group is
\[
 \Cl_{(6)\infty_2}(K_6)\cong C_6.
\tag{28}
\]
Its order-two character gives the lower-conductor stratum.  If
\[
 Y=\frac{\beta_6-2+\sqrt{\beta_6-1}}2>1,
\]
then the modulus-three audit
\path{scripts/dimension_six_lower_stratum.gp}, combining the analytic
class-number formula with Kopp's Theorem~1.1~\cite{Kopp}, proves
unconditionally that the overlap $y=\nu_{(0,2)}$ satisfies
\[
 y^2=Y,\qquad y^2+y^{-2}=\beta_6-2.
\tag{29}
\]

The remaining three reciprocal pairs form the primitive order-six
packet.  Let
\[
\begin{aligned}
 \mathcal Q_6(X)={}&X^{12}+3X^{11}-6X^{10}-16X^9+3X^8+27X^6\\
 &\quad+3X^4-16X^3-6X^2+3X+1,
\end{aligned}
\tag{30}
\]
and let $x_{\rm alg}$ be its unique root in
\[
 2.21288528901718<x_{\rm alg}<2.21288528901719.
\tag{31}
\]
The norm-$37$ prime $(4\beta_6+1)$ generates the ray group and fixes the
arithmetic-Frobenius order of the other two primitive pairs.  Inserting
this Artin-labeled algebraic packet together with (29) in the
six-dimensional Weyl reconstruction makes all
\[
 \binom62^2=225
\]
two-by-two minors vanish exactly; the reconstructed matrix has trace one.
Thus the finite algebra has no remaining ambiguity.

In the double-sine convention (1), the distinguished
$(0,1)$ analytic value is
\[
 x_{\rm an}
 =
 \frac{
 S_2(\beta_6/6\mid\beta_6,1)
 S_2(1/6\mid\beta_6,1)}
 {S_2((\beta_6+1)/6\mid\beta_6,1)}.
\tag{32}
\]

\paragraph{Exact closure target.}
The exact finite audit shows the following conditional implication:
the canonical analytic table agrees with the certified rank-one
algebraic table, and hence gives both formal shifts, if
\[
 x_{\rm an}=x_{\rm alg}.
\tag{33}
\]
The characteristic orbits, reciprocity, and the arithmetic-Frobenius
labeling propagate this one equality through the other primitive
pairs.  The script
\path{scripts/dimension_six_primitive_fourier_audit.gp} records the
equivalent oriented order-six Fourier identity, including its ray-group
generator, character convention, Euler-factor convention, and
embedding order.  That analytic identity, and therefore (33), remains
unproved here.

This is not a finite-computation issue.  The ray field has signature
$(6,3)$; its primitive character has order six; and its cubic component
is wildly ramified above $3$.  Shintani's index-two theorem does not
apply, Roblot's hypotheses~\cite{Roblot2013} do not cover this wild
sextic case, and
proved CM Brumer--Stark results do not apply to this mixed-signature
extension.  Standard double-sine shifts, reflection, duplication, and
conductor lowering determine norms of the four level-six lifts but do
not isolate the selected lift.

There is also no weaker local-divisibility shortcut.  For
\[
 A_6=\begin{pmatrix}115&-24\\24&-5\end{pmatrix},
\]
write $A_6\cdot z=(115z-24)/(24z-5)$.  Then
\[
 A_6\cdot z-z
 =-24\frac{z^2-5z+1}{24z-5}.
\tag{34}
\]
The right side has a simple zero at $\beta_6$.  Consequently, in the
local holomorphic ring at $\beta_6$, divisibility of a defect entry by
$A_6\cdot z-z$ is equivalent to its vanishing at $\beta_6$; proving
such divisibility would already prove the missing TCC identity.

\subsection{Dimension eight: an unconditional lower stratum and two
quartic phases}

Let
\[
 K_8=\Q(\sqrt5),\qquad
 \phi=\frac{1+\sqrt5}{2},\qquad
 \beta_8=3\phi+2=\frac{7+3\sqrt5}{2}.
\]
The canonical form has discriminant $45$ and belongs to the order
\[
 \mathcal O_3=\Z+3\mathcal O_{K_8}.
\]
The order-ray group and its maximal-order realization are
\[
 \Cl_{(8)\infty_2}(\mathcal O_3)
 \cong
 \Cl_{(24)\infty_2}(\mathcal O_{K_8})
 \cong C_4\times C_2\times C_2.
\tag{35}
\]
The change-of-order isomorphism in (35) does not identify partial zeta
functions by itself.  Kopp's conductor-lowering formula~\cite{Kopp}
supplies the
required three-factor bridge between the order-modulus-eight
characteristics and maximal-order moduli $8$ and $24$.

The fifteen nonzero lower-conductor characteristics lie in the
one-place ray-$12$ field.  Its one-place and both-place ray groups are
$C_2^2$ and $C_2^3$, and the extension is in Shintani's index-two
range.  An exact divisor audit gives safe exponent $576$.  Arb
enclosures and the height gap used above then prove unconditionally
that its four independent positive values are the real roots of
\[
\begin{aligned}
 P_{\rm low}(X)={}&X^8-8X^7+12X^6+8X^5-22X^4\\
                  &\quad+8X^3+12X^2-8X+1.
\end{aligned}
\tag{36}
\]
Zauner covariance and reciprocity supply the full lower stratum.

The forty-eight primitive characteristics contain four quadratic
Fourier components and two conjugate pairs of quartic components.
The two cyclic quartic quotient fields have relative equations
\[
\begin{aligned}
 F_0/K_8:\quad&X^4-6\phi X^2+3,\\
 F_1/K_8:\quad&X^4+(6-6\phi)X^2+3.
\end{aligned}
\tag{37}
\]
The exact models, fundamental-unit coordinates, and embeddings used
below are fixed in
\path{scripts/dimension_eight_quartic_bridge.gp}.  Let
$\eta_0,\eta_1$ be the anti-units selected there and let
$\rho_1,\ldots,\rho_4$ denote the four real embeddings in its
local-Frobenius order.  Define the oriented logarithmic resolvents
\[
\begin{aligned}
 R_0&=\log|\eta_0|_{\rho_3}
      +i\log|\eta_0|_{\rho_4},\\
 R_1&=\log|\eta_1|_{\rho_2}
      +i\log|\eta_1|_{\rho_1}.
\end{aligned}
\tag{38}
\]
In the PARI dual coordinates for the cyclic basis
$C_4\times C_2\times C_2$ fixed by that script, Roblot's
cyclic-quartic theorem proves the two modulus identities
~\cite{Roblot2013}
\[
 |L'_S(0,[1,0,0])|=|R_0|,
 \qquad
 |L'_S(0,[1,1,0])|=|R_1|.
\tag{39}
\]
It does not prove the orientations.

\paragraph{Projective CM-descent gate.}
There is nevertheless a sharper route than an abstract oriented
refinement of (39).  To handle the infinite-place swap exactly, embed
the one-place ray group into the full group
\[
 \Cl_{(24)\infty_1\infty_2}(K_8)\cong C_4\times C_2^3.
\]
In its PARI cyclic coordinates, base conjugation has matrix
\[
 C_\sigma=
 \begin{pmatrix}
 3&2&0&2\\0&1&0&0\\0&1&1&0\\0&0&0&1
 \end{pmatrix}.
\tag{CM1}
\]
The pullbacks of the two one-place characters, their conjugates, and
their inverses have full-ray dual coordinates
\[
\begin{array}{c|c|c|c|c}
\chi&\chi&\chi^\sigma&\chi^{-1}
 &\chi(\chi^\sigma)^{-1}\\ \hline
{[1,0,0]}&[3,0,0,0]&[1,1,0,1]&[1,0,0,0]&[2,1,0,1]\\
{[1,1,0]}&[3,0,0,1]&[1,1,0,0]&[1,0,0,1]&[2,1,0,1].
\end{array}
\tag{CM2}
\]
Thus the literal identity $\chi^\sigma=\chi^{-1}$ fails: its two sides
have different labeled infinite-place conductors.  However, the
projectively decisive quotient is the same order-two character in both
rows.  Its conductor is $(24)\infty_1\infty_2$, and exact class-field
construction gives
\[
 K_8(\sqrt{-6})/\Q,\qquad
 \{\text{quadratic subfields}\}
 =\{\Q(\sqrt5),\Q(\sqrt{-6}),\Q(\sqrt{-30})\}.
\tag{CM3}
\]
Consequently the induced Artin representations have projective image
$V_4$, not a larger dihedral group.  Equations (CM1)--(CM3), including
the one-place quotient maps, are reproduced by
\path{scripts/dimension_eight_cm_descent.gp}.

The exact audit script
\path{scripts/dimension_eight_exact_tcc.gp} constructs an algebraic
candidate packet: its primitive squared-overlap field has degree $32$
over $\Q$, all signed overlaps lie in one degree-$64$ field, and
adjoining the Weyl phases gives a convention-selected degree-$128$
compositum.  For each of the two formal shifts, the script verifies
\[
 \Tr\Pi=1,\qquad \Pi^2=\Pi,\qquad
 \text{all }\binom82^2=784\text{ two-by-two minors vanish}.
\tag{40}
\]

\paragraph{Exact closure targets.}
The canonical analytic table would coincide with this candidate packet
and therefore satisfy both formal shifts if the two convention-matched
oriented identities
\[
\begin{aligned}
 L'_S(0,[1,0,0])&=R_0,\\
 L'_S(0,[1,1,0])&=R_1
\end{aligned}
\tag{41}
\]
hold.  The lower-conductor special values and the algebraic
calculation \textup{(40)} is unconditional; its identification with
the primitive analytic table is conditional on \textup{(41)}.

Rigorous double-sine windows locate all eight independent primitive
analytic squares in the corresponding isolated algebraic intervals.
At high precision the phase quotients in (41) equal
$1+O(10^{-114})$.  This is strong evidence, not an equality proof:
(39) permits an arbitrary unit-circle phase.  The exact finite TCC
orientation sieve selects one pair among the $64$ natural algebraic
orientations, but using that selection to deduce (41) would be circular.
The projective computation above narrows the missing proof further.
One must construct the two re-induced Hecke characters over
$\Q(\sqrt{-6})$ or $\Q(\sqrt{-30})$, verify the linear (not merely
projective) Artin induction including any scalar twist, match the
removed Euler factors defining $L_S$, and translate the
imaginary-quadratic Kronecker-limit formula into the orientations
$R_0,R_1$.  That CM calculation would supply the phase and make (41),
and hence dimension eight, unconditional without further finite
computation.  The present paper records the exact descent gate but
does not assert that final convention match.

\subsection{Comparison}

The two composite dimensions fail for different reasons:
\[
\begin{array}{c|c|c|c}
d&\text{unconditional analytic part}&\text{exact finite part}
 &\text{missing input}\\ \hline
6&\text{lower stratum, modulus and sign}
 &225\text{ minors}
 &1\text{ order-six value}\\
8&15\text{ lower overlaps, quartic moduli, projective CM gate}
 &2\times784\text{ minors}
 &\text{linear CM re-induction and two phase matches}.
\end{array}
\tag{42}
\]
Dimension six has no analogous shortcut: its primitive character has
order six, so the corresponding projective quotient has order three
rather than two.  Dimension eight is therefore closer to closure than
the raw comparison (39) suggests; its remaining obstruction is the
linear re-induction (including a possible scalar twist), the
$S$-Euler-factor comparison, and the two imaginary-quadratic
Kronecker-limit convention matches.

\section{Scope and reproducibility}

Theorem~\ref{thm:tcc} proves the complete formal TCC in dimension five.
Together with Section~\ref{sec:d4}, it proves
Theorem~\ref{thm:lowdim}.  It does not prove a dimension-five SIC,
the TCC in arbitrary
dimensions, or Zauner's conjecture.  The new number-theoretic step is
Theorem~\ref{thm:packet}: it proves the previously conjectural
dimension-five ray packet by combining Kopp's proved limit formula with
the applicable case of Shintani's proved algebraicity theorem and an
effective height argument.

Section~\ref{sec:composite} records research-status closure targets,
not conditional propositions masquerading as higher-dimensional
theorems.  It shows how far the proved low-dimensional machinery
extends: the finite algebraic candidate packets pass exact cancellation
audits, while their identification with the primitive analytic values
still requires (33) in dimension six and the two orientations (41) in
dimension eight.  More decimal recognition cannot replace those
missing special-value proofs.

All reported certificates were regenerated on Linux
6.8.0-124-generic (x86-64, four virtual AMD EPYC 9354P cores) with
PARI/GP~2.15.4, Python~3.12.3, python-flint~0.9.0, and FLINT~3.6.0.
Representative wall times and peak resident memory were:
\[
\begin{array}{lrr}
\text{audit}&\text{wall time}&\text{peak RSS}\\ \hline
\text{Shintani field and exponent}&0.9\ {\rm s}&22\ {\rm MB}\\
\text{dimension-five exact minors}&1.1\ {\rm s}&22\ {\rm MB}\\
\text{embedding/factor selection}&0.3\ {\rm s}&23\ {\rm MB}\\
\text{Arb double-sine enclosures}&33.8\ {\rm s}&30\ {\rm MB}.
\end{array}
\]
The complete regression suite takes approximately three minutes on
that platform.  A tagged immutable source archive, including the files
below and \texttt{certificates/SHA256SUMS}, will be deposited with
Zenodo and its DOI inserted in the final submission.  A repository
commit hash alone is not intended to substitute for that persistent
archive.

The supplementary package contains:
\begin{itemize}
\item the dimension-four $36$-minor, ray, multiplier, and PARI
  certificates;
\item the dimension-five 25-characteristic bridge and exact
  $C_8$ Fourier-support audit;
\item the PARI transcript with \texttt{bnfcertify}=1, positive
  ray-class logs, and Sturm root isolation;
\item the exact interval/conjugate/factor embedding certificate,
  formal 100-minor Laurent certificate, and four-factor reduction;
\item the Shintani sign-class, fixed-field, conductor, and
  divisor-by-divisor exponent-$5760$ audit;
\item the labeled $K$-isomorphism, unit, Sturm, Frobenius-orbit,
  and rigorous Arb height-gap certificates;
\item the dimension-six ray, lower-stratum, Frobenius-orientation,
  $225$-minor, and fixed-point-factorization audits; and
\item the dimension-eight conductor-lowering, lower-stratum,
  cyclic-quartic, projective CM-descent, Artin-label, orientation,
  and two-shift
  $784$-minor audits.
\end{itemize}
SHA-256 digests for the archived dimension-four and dimension-five
artifacts are recorded in \texttt{certificates/SHA256SUMS}; the
dimension-six and dimension-eight research scripts are fixed by the
repository commit accompanying this manuscript.

\section*{Acknowledgments}

The author acknowledges the use of OpenAI's GPT-5.6 Sol model for
mathematical brainstorming, symbolic-computation scaffolding, code
review, and manuscript editing.  All proofs, computations, and claims
were independently checked through the exact certificates described
above, and responsibility for the final manuscript rests with the
author.  Additional personal and institutional acknowledgments will be
added before submission.

\begin{thebibliography}{12}
\bibitem{Kopp}
G.~Kopp, \emph{The Shintani--Faddeev modular cocycle: Stark units from
$q$-Pochhammer ratios}, arXiv:2411.06763.
\bibitem{AFK}
D.~M.~Appleby, S.~T.~Flammia, and G.~S.~Kopp,
\emph{A constructive approach to Zauner's conjecture via the Stark
conjectures}, arXiv:2501.03970.
\bibitem{Kopp2018}
G.~Kopp, \emph{SIC-POVMs and the Stark conjectures},
Int.\ Math.\ Res.\ Not.\ IMRN 2021, no.~18, 13812--13838;
arXiv:1807.05877.
\bibitem{Kopp2021}
G.~S.~Kopp, \emph{Indefinite zeta functions},
Res.\ Math.\ Sci.\ \textbf{8}, 17 (2021),
\url{https://doi.org/10.1007/s40687-021-00252-9}.
\bibitem{Kopp2023}
G.~S.~Kopp, \emph{A Kronecker limit formula for indefinite zeta
functions}, Res.\ Math.\ Sci.\ \textbf{10}, 24 (2023),
\url{https://doi.org/10.1007/s40687-023-00384-0}.
\bibitem{Shintani1978}
T.~Shintani, \emph{On certain ray class invariants of real quadratic
fields}, J.\ Math.\ Soc.\ Japan \textbf{30} (1978), 139--167,
\url{https://doi.org/10.2969/jmsj/03010139}.
\bibitem{Voutier1996}
P.~Voutier, \emph{An effective lower bound for the height of algebraic
numbers}, Acta Arith.\ \textbf{74} (1996), 81--95,
\url{https://doi.org/10.4064/aa-74-1-81-95}.
\bibitem{PARI}
The PARI Group, \emph{PARI/GP version 2.15.4},
Université de Bordeaux, 2023,
\url{https://pari.math.u-bordeaux.fr/}.
\bibitem{pythonflint}
F.~Johansson and contributors, \emph{python-flint 0.9.0},
\url{https://github.com/flintlib/python-flint}.
\bibitem{Roblot2013}
X.-F.~Roblot, \emph{Index formulae for Stark units and their
solutions}, Pacific J.\ Math.\ \textbf{266} (2013), 391--422,
\url{https://doi.org/10.2140/pjm.2013.266.391}.
\end{thebibliography}

\end{document}
